% TOP_PROBLEM: {"id": 30006693, "title": "Grothendieck's Standard Conjecture D", "category_id": 5, "category": {"id": 5, "name": "algebraic_geometry", "display_name": "Algebraic Geometry", "description": "Geometric objects defined by polynomial equations.", "slug": "algebraic-geometry", "order_index": 5, "created_at": "2026-07-31T15:26:25.671Z"}, "status": "open"}
% FIRST_POSTED: 2026-10-06
% ATTEMPT_STATUS: unresolved
% OPENAI_CLAIM: {"schema_version":1,"source_repo":"https://github.com/openai/math","source_commit":"adc7f1241b42e322a6451854ab7e4b4c146bf78a","release_date":"2026-10-06","match":"partial","resolution":"partial","independently_reviewed":false,"second_pass":"n/a","families":[{"family":"032","title":"Hodge and Kuga–Satake results for all projective K3 surfaces","subject":"Algebraic and complex geometry","manuscripts":[{"title":"The rational Hodge conjecture for CM abelian varieties","pdf_path":"preprints/The-rational-Hodge-conjecture-for-CM-abelian-varieties-September-30-2026/paper.pdf","readme_path":"preprints/The-rational-Hodge-conjecture-for-CM-abelian-varieties-September-30-2026/README.md","date":"2026-09-30","theorem_ref":"build/sections/07-consequences.tex:274, Corollary cc:standard"}],"lean":null,"reasoning_trace":null}]}
\documentclass[11pt]{article}
\usepackage[margin=25mm]{geometry}
\usepackage{fontspec}
\setmainfont{Latin Modern Roman}
\usepackage{amsmath,amssymb,mathtools}
\usepackage{unicode-math}
\setmathfont{Latin Modern Math}
\usepackage{enumitem}
\usepackage{xurl,hyperref,bookmark}
\hypersetup{colorlinks=true,urlcolor=blue}
\setcounter{secnumdepth}{0}
\setlength{\parindent}{0pt}
\setlength{\parskip}{5pt}
\setlength{\emergencystretch}{3em}
\begin{document}
\section*{Grothendieck's Standard Conjecture D}

\subsection{OpenAI's claim}
The manuscripts claim algebraicity of rational Hodge classes on CM abelian varieties, with consequences for generalized Hodge classes and finite-field Tate classes. The cited manuscript argues for that result; its relation to this problem is specified below.

\subsection{Outline of the argument}
Manuscript 1 (family 032, 2026-09-30): The argument realizes four-factor balanced tensors using unitary theta constructions. Elementary switches between equal-sum lists, followed by algebraic polarization contractions and descent of coefficients, build arbitrary balanced CM tensors from these elementary cycle classes.

\subsection{Scope relative to this problem}
Manuscript 1 (family 032, 2026-09-30): The standard corollary identifies numerical and prime-to-characteristic etale homological equivalence for abelian varieties; our target ranges over all smooth projective varieties and available Weil cohomology theories.

\subsection{Formal verification}
Family 032 has no Lean scope page in this snapshot; no formalization of these manuscript results is listed.

This site has not executed the Lean files or checked their dependencies. These descriptions report the scope of the linked declarations, rather than a new verification of the proofs.

\subsection{Status caveat}
OpenAI attributes these results to an internal model. Its release notice cautions that results lacking full formalization can contain mistakes. This site has not independently assessed the proofs; a solved label records the scope of the reported claim.

\subsection{Sources}
\begin{itemize}
\item OpenAI, \emph{The rational Hodge conjecture for CM abelian varieties}, 2026-09-30; citation key \texttt{OAI:The-rational-Hodge-conjecture-for-CM-abelian-varieties-September-30-2026}. \href{https://github.com/openai/math/blob/adc7f1241b42e322a6451854ab7e4b4c146bf78a/preprints/The-rational-Hodge-conjecture-for-CM-abelian-varieties-September-30-2026/paper.pdf}{PDF (GitHub)}; \href{https://github.com/openai/math/raw/adc7f1241b42e322a6451854ab7e4b4c146bf78a/preprints/The-rational-Hodge-conjecture-for-CM-abelian-varieties-September-30-2026/paper.pdf}{PDF (direct)}; \href{https://github.com/openai/math/blob/adc7f1241b42e322a6451854ab7e4b4c146bf78a/preprints/The-rational-Hodge-conjecture-for-CM-abelian-varieties-September-30-2026/README.md}{Manuscript page and citation}; \href{https://github.com/openai/math/blob/main/preprints/The-rational-Hodge-conjecture-for-CM-abelian-varieties-September-30-2026/paper.pdf}{Latest version}.
\item \href{https://github.com/openai/math/blob/adc7f1241b42e322a6451854ab7e4b4c146bf78a/CONTENTS.md}{OpenAI catalogue entry}.
\end{itemize}
\end{document}
